% GENERATED FILE. Edit canonical lessons, then run npm run generate:books.
% canonical-source-hash: fnv1a64:ca90f41d17010d1f
% canonical-lessons: FR-C494-elever, FR-C494-gater, FR-C494-punition, FR-C494-desobeir, FR-C494-creche

\chapter{Élever des enfants}
\label{ch:fr-elever-des-enfants}
\hlchaptermodality{494}

\begin{chapteropening}
\textbf{By the end of this chapter:} \emph{I can talk about bringing up children.}

Complete the last lesson of chapter 494: I can talk about bringing up children.
\end{chapteropening}

\section[élever]{élever — to bring up, to raise}

\label{lesson:FR-C494-elever}



\begin{quote}
Before the new one: say the French for baby's bottle, then the French for to rock.
\end{quote}

\subsection*{You'll want to know: élever}
\textbf{élever} — \textquotedblleft{}to bring up, to raise\textquotedblright{}.

\begin{quote}
Ses grands-parents l'ont élevé à la campagne.
\end{quote}

\begin{quote}
\emph{His grandparents brought him up in the country.}
\end{quote}

Accent change: \emph{j'élève}, \emph{nous élevons}. Also to raise animals: \emph{élever des moutons}.

\begin{cousinweb}
\emph{élever} is built from \emph{é-} and \emph{lever} \textquotedblleft{}to raise\textquotedblright{}, which is inherited from Latin \emph{levāre} \textquotedblleft{}to lift\textquotedblright{} (Lewis \& Short).
\end{cousinweb}

\begin{grammarlens}[title={What you've built}]
One more word for telling, explaining and describing at B1.
\end{grammarlens}

\subsection*{Your turn}
\begin{itemize}
\raggedright
  \item \emph{Say it:} \emph{élever}
  \item \emph{Say it:} \emph{élever}, once more
  \item \emph{Say it:} a sentence of your own with \emph{élever}
  \item \emph{Recall:} say the French for baby's bottle, then the French for to rock, then say \emph{élever} again
\end{itemize}

\begin{tcolorbox}[breakable,colback=teal!4,colframe=teal!35!black,title={Before you move on}]
What does élever mean? (\textquotedblleft{}To bring up, to raise\textquotedblright{}.) Say it once more.
\end{tcolorbox}
\section[gâter]{gâter — to spoil (a child); to ruin}

\label{lesson:FR-C494-gater}



\begin{quote}
Before the new one: say the French for to rock, then the French for to bring up.
\end{quote}

\subsection*{You'll want to know: gâter}
\textbf{gâter} — \textquotedblleft{}to spoil (a child); to ruin\textquotedblright{}.

\begin{quote}
Mon oncle gâte beaucoup ses neveux.
\end{quote}

\begin{quote}
\emph{My uncle really spoils his nephews.}
\end{quote}

\emph{Un enfant gâté} — a spoilt child. \emph{Il fait un temps gâté} — the weather has turned bad.

\begin{cousinweb}
\emph{gâter} is inherited from Latin \emph{vastāre} \textquotedblleft{}to lay waste\textquotedblright{} (EtymDB; Lewis \& Short); its first sound was reshaped under a related Germanic word.
\end{cousinweb}

\begin{grammarlens}[title={What you've built}]
One more word for telling, explaining and describing at B1.
\end{grammarlens}

\subsection*{Your turn}
\begin{itemize}
\raggedright
  \item \emph{Say it:} \emph{gâter}
  \item \emph{Say it:} \emph{gâter}, once more
  \item \emph{Say it:} a sentence of your own with \emph{gâter}
  \item \emph{Recall:} say the French for to rock, then the French for to bring up, then say \emph{gâter} again
\end{itemize}

\begin{tcolorbox}[breakable,colback=teal!4,colframe=teal!35!black,title={Before you move on}]
What does gâter mean? (\textquotedblleft{}To spoil (a child); to ruin\textquotedblright{}.) Say it once more.
\end{tcolorbox}
\section[la punition]{la punition — punishment}

\label{lesson:FR-C494-punition}



\begin{quote}
Before the new one: say the French for to bring up, then the French for to spoil.
\end{quote}

\subsection*{You'll want to know: la punition}
\textbf{la punition} — \textquotedblleft{}punishment\textquotedblright{}.

\begin{quote}
Comme punition, il n'a pas eu de dessert.
\end{quote}

\begin{quote}
\emph{As a punishment, he did not get any pudding.}
\end{quote}

The verb is \emph{punir} \textquotedblleft{}to punish\textquotedblright{}.

\begin{cousinweb}
\emph{la punition} is a learned borrowing from Latin \emph{pūnītiō} \textquotedblleft{}punishment\textquotedblright{} (EtymDB), from \emph{pūnīre} \textquotedblleft{}to punish\textquotedblright{}, built on \emph{poena} \textquotedblleft{}penalty\textquotedblright{} (Lewis \& Short).
\end{cousinweb}

\begin{grammarlens}[title={What you've built}]
One more word for telling, explaining and describing at B1.
\end{grammarlens}

\subsection*{Your turn}
\begin{itemize}
\raggedright
  \item \emph{Say it:} \emph{la punition}
  \item \emph{Say it:} \emph{la punition}, once more
  \item \emph{Say it:} a sentence of your own with \emph{la punition}
  \item \emph{Recall:} say the French for to bring up, then the French for to spoil, then say \emph{la punition} again
\end{itemize}

\begin{tcolorbox}[breakable,colback=teal!4,colframe=teal!35!black,title={Before you move on}]
What does la punition mean? (\textquotedblleft{}Punishment\textquotedblright{}.) Say it once more.
\end{tcolorbox}
\section[désobéir]{désobéir — to disobey}

\label{lesson:FR-C494-desobeir}



\begin{quote}
Before the new one: say the French for to spoil, then the French for punishment.
\end{quote}

\subsection*{You'll want to know: désobéir}
\textbf{désobéir} — \textquotedblleft{}to disobey\textquotedblright{}.

\begin{quote}
Le chien désobéit quand il voit un chat.
\end{quote}

\begin{quote}
\emph{The dog disobeys when it sees a cat.}
\end{quote}

\emph{Désobéir à quelqu'un}. A regular \emph{-ir} verb: \emph{ils désobéissent}.

\begin{cousinweb}
\emph{désobéir} is built from \emph{dés-} and \emph{obéir} \textquotedblleft{}to obey\textquotedblright{}, which goes back to Latin \emph{oboedīre} \textquotedblleft{}to listen to, to obey\textquotedblright{}, built on \emph{audīre} \textquotedblleft{}to hear\textquotedblright{} (Lewis \& Short).
\end{cousinweb}

\begin{grammarlens}[title={What you've built}]
One more word for telling, explaining and describing at B1.
\end{grammarlens}

\subsection*{Your turn}
\begin{itemize}
\raggedright
  \item \emph{Say it:} \emph{désobéir}
  \item \emph{Say it:} \emph{désobéir}, once more
  \item \emph{Say it:} a sentence of your own with \emph{désobéir}
  \item \emph{Recall:} say the French for to spoil, then the French for punishment, then say \emph{désobéir} again
\end{itemize}

\begin{tcolorbox}[breakable,colback=teal!4,colframe=teal!35!black,title={Before you move on}]
What does désobéir mean? (\textquotedblleft{}To disobey\textquotedblright{}.) Say it once more.
\end{tcolorbox}
\section[la crèche]{la crèche — day nursery, crèche; nativity scene}

\label{lesson:FR-C494-creche}



\begin{quote}
Before the new one: say the French for punishment, then the French for to disobey.
\end{quote}

\subsection*{You'll want to know: la crèche}
\textbf{la crèche} — \textquotedblleft{}day nursery, crèche; nativity scene\textquotedblright{}.

\begin{quote}
Le matin, je dépose ma fille à la crèche.
\end{quote}

\begin{quote}
\emph{In the morning I drop my daughter off at the nursery.}
\end{quote}

For children under three. In Belgium also \emph{la crèche}; in Québec \emph{la garderie}.

\begin{cousinweb}
\emph{la crèche} comes from Frankish \emph{krippija} \textquotedblleft{}a manger, a crib\textquotedblright{} (EtymDB), a Germanic word related to English \emph{crib}.
\end{cousinweb}

\begin{grammarlens}[title={What you've built}]
One more word for telling, explaining and describing at B1.
\end{grammarlens}

\subsection*{Your turn}
\begin{itemize}
\raggedright
  \item \emph{Say it:} \emph{la crèche}
  \item \emph{Say it:} \emph{la crèche}, once more
  \item \emph{Say it:} a sentence of your own with \emph{la crèche}
  \item \emph{Recall:} say the French for punishment, then the French for to disobey, then say \emph{la crèche} again
\end{itemize}

\begin{tcolorbox}[breakable,colback=teal!4,colframe=teal!35!black,title={Before you move on}]
What does la crèche mean? (\textquotedblleft{}Day nursery, crèche; nativity scene\textquotedblright{}.) Say it once more.
\end{tcolorbox}
